\documentclass[10pt,a4paper]{amsart}
\usepackage[margin=19mm]{geometry}
\usepackage{amsmath,amssymb,mathtools}
\usepackage[scr=rsfs,cal=euler]{mathalfa}
\usepackage{xfrac}
\usepackage[unicode,hidelinks,bookmarks=false]{hyperref}
\newcommand{\ie}{i.e.\ }
\newcommand{\cf}{cf.\ }
\newcommand{\resp}{respectively\ }
\newcommand{\Subsection}{\subsection}
\providecommand{\nobreakdash}{\nobreak\hbox{-}}
\setlength{\emergencystretch}{2em}
\multlinegap0pt
\numberwithin{equation}{section}
\newtheorem{theorem}{Theorem}[section]
\newtheorem{lemma}[theorem]{Lemma}
\newtheorem{corollary}[theorem]{Corollary}
\newtheorem{remark}[theorem]{Remark}
\allowdisplaybreaks
\title[Stationary Navier-Stokes with nonlinear damping: original Theorem 1.2]{Variational approach for the $n$-dimensional stationary Navier-Stokes equations with a damping term: original Theorem 1.2 and its complete original proof}
\author{Alireza Khatib, Abbas Moameni and Somayeh Mousavinasr}
\date{Source: Electronic Journal of Differential Equations 2026 (2026), No. 01, pp. 1--13; DOI 10.58997/ejde.2026.01}
\begin{document}
\maketitle
\noindent\textit{Editorial scope.} The counted unit of this excerpt is the article's own Theorem 1.2 as printed, together with its complete original proof. Theorem 1.1 (the paper's minimax principle) with its own complete proof, and the whole of the original Section 2 --- weak closedness of $K(r)$, Ky Fan's minimax principle, the weak continuity and weak lower semicontinuity lemmas, the existence of $\bar u$, and the estimates, Stokes and elliptic lemmas used by the proof --- are retained verbatim as uncounted context, so that the counted proof can be read without leaving the excerpt. The article's Theorems 1.3 and 1.4, Remark 1.5, the symmetry discussion and Section 3 are outside this unit. Printed slips are kept literal and no mathematics has been rewritten or repaired: the source prints \textit{Fro each} in Lemma 2.4, prints \textit{rVert} (a missing backslash in the source) in Lemmas 2.8 and 2.9, prints a lowercase exponent in $w^{2,2}$ in (2.12), and writes $\bar u^{p-1}$ for $g(x)$ in the proof where Lemma 2.6 uses $\bar u^{p-2}$. The source labels the estimate (1.3) twice, in Theorem 1.2 and again inside the proof; the redundant duplicate label command in the proof is the only source string not carried over (both displayed formulas are kept verbatim and no reference uses that label).
\section{Introduction}

In this article we study  the $n$-dimensional  stationary Navier-Stokes equations
with a damping term,
\begin{equation}\label{eq1}
 \begin{gathered}
 -\Delta u+ (u \cdot \nabla)u  +\mu |u|^{p-2}u = f(x) + \nabla P  \quad
   \forall x\in \Omega\\
 \nabla  u=0  \quad  \forall x\in \Omega \\
  u=0  \quad   \forall x\in \partial\Omega
\end{gathered}
\end{equation}
where \(\Omega \subset \mathbb{R} ^n\) is  bounded,  $p\geq 1$ and
$\mu \in \mathbb{R}$.  We address both  linear and nonlinear dampings
and we are allowing $\mu$ to take both positive and negative values.
Here \(u  = (u_1, u_2,\dots, u_n)\) is the velocity, \(P\) stand for scalar pressure
and \(f\) is the  external force.

The analysis of the Navier-Stokes equations is a central theme in mathematical fluid
mechanics. For the classical evolutionary system without damping, the existence of
global weak solutions was established by Leray \cite{JL} and Hopf \cite{EH}, but the
uniqueness of weak solutions and the global existence of strong solutions remain open
problems. These longstanding challenges have motivated researchers to study modified
models where additional terms improve the mathematical structure of the equations.

One such modification is the inclusion of a damping term of the form $\mu |u|^{p-2}u$.
From the physical viewpoint, this term models resistance to motion and arises naturally
in contexts such as porous media flows, drag or friction effects, and other dissipative
mechanisms (see \cite{LHS, RTE}). From the mathematical viewpoint, damping often
leads to better control of solutions, sometimes yielding results that are out of reach
for the standard Navier-Stokes equations. This has stimulated extensive research on the
evolutionary Navier-Stokes system with damping, leading to results on global existence,
regularity, decay rates, attractors, and stability (see, e.g.,
\cite{XC, YJX, ZHJ, ZHM, XLS,ZJZ, YZ, ZZC}).

In contrast,  stationary Navier-Stokes equations with damping have received comparatively
less attention. Some existence and uniqueness results are available for the case
$\mu > 0$ \cite{WLW}, and there is growing literature on numerical approaches
(see \cite{ZLD, LML, HQL, HQY, ZBS}). However, general variational approaches capable
of handling both positive and negative damping, as well as borderline cases where
coercivity is lost, are still lacking.

In this work, we develop a new minimax principle on convex subsets of Banach spaces and
apply it to the $n$-dimensional stationary Navier-Stokes equations with damping.
Our framework is broad and flexible: depending on the choice of convex set, it allows
us to obtain solutions with additional structural properties, such as symmetry.
With this method we prove existence results for both linear ($p=2$) and nonlinear
damping, without restriction on the damping constant, and even in the critical case
$\mu = -\lambda_1$. This general minimax principle is of independent interest and may
find applications to other nonlinear PDEs beyond the Navier-Stokes system.

In this work,  we  consider the Banach space
\[
V= \{ u\in H_0^1(\Omega)\cap L^p(\Omega),\nabla . u=0\},
\]
equipped with the  norm
\[
\rVert u \rVert:= \rVert u \rVert_{ H_0^{1}(\Omega)}
+\rVert u \rVert_{ L^{p}(\Omega)}.
\]
Let \(\Lambda u\) be the operator \(\Lambda u := (u . \nabla)u\),
and \( K\) be a convex and weakly closed subset of \(V\).
We shall define \( M : K \times K \to\mathbb{R}\) as follows,
\begin{equation}\label{def II}
    M(u,v)=\frac{1}{2}\int_{\Omega} |\nabla u|^2\,dx
    -\frac{1}{2}\int_{\Omega} |\nabla v|^2\,dx
    +\int_\Omega (\Lambda u-f(x)  -\frac{1}{p}|u|^{p-2}u) (u-v)\,dx,
\end{equation}
where \(f \in L^2(\Omega)\). The following variational principle on general convex
sets $K$  is a  key component in our arguments. It is also broad enough to deal with
various other cases by choosing a convex set $K$ accordingly.

\setcounter{theorem}{0}
\begin{theorem}\label{variational}
Let  \( K\) be a convex and weakly closed subset of \(V\). Assume that the following
two assertions hold:
\begin{itemize}
\item[(i)] There  exists \(\bar{u} \in K\) for   such that
    \[
    M(\bar{u}, v)\leq 0,\quad  \forall v\in  K,
    \]
    where \(M \) is defined in \eqref{def II}.

\item[(ii)] There exists \(\bar{v} \in K\) such that \[-\Delta \bar{v} + \nabla P = f(x) +|\bar{u}|^{p-2}\bar{u} - \Lambda \bar{u},\]
    in the weak sense, i.e.,
\[
   \int_\Omega\nabla\bar{v}.\nabla\eta\,dx
   =\int_\Omega (f(x)  +|\bar{u}|^{p-2}\bar{u} -
   \Lambda \bar{u})\eta\,dx,\quad \forall\eta\in V.
\]
\end{itemize}
Then \(\bar{u} \in K\) is a weak solution of the equation
\[
-\Delta u+ \;\Lambda u \; = f(x) + \nabla P +|u|^{p-2}u \,.
\]
\end{theorem}

\setcounter{theorem}{1}
\begin{theorem}\label{main theorem 1}
Let \(\Omega\) be a  bounded \(C^2\) domain in \(\mathbb{R} ^n\) and  \(\mu<0\).
Then for \(f\in L^2(\Omega) \) small enough,  the following statements hold:
 \begin{itemize}
 \item[(i)] For \(n\leq 4\) and \(p>2\), the Navier-Stokes equation
     \eqref{eq1} has a  solution \(u \in  W^{2,2}(\Omega)  \).

 \item[(ii)] For \(5\leq n\leq 7\) and  \( 2< p \leq \frac{2n-4}{n-4}\),
 the Navier-Stokes equation \eqref{eq1} has a  solution in \( W^{2,2}(\Omega)\).
 \end{itemize}
In both cases there exists a scalar function $P:\Omega \to \mathbb{R}$ and a constant $C>0$
such that
\begin{equation} \label{nabla11}
     \rVert \Delta {u} \rVert_{L^2(\Omega)}+\rVert \nabla P \rVert_{L^2(\Omega)}
     \leq C\big( \rVert f\rVert_{L^2(\Omega)}+\rVert u\rVert^{p-1}_{L^{2(p-1)}}
     +\rVert {u}\rVert^{2}_{W^{1,2^*}(\Omega)} \rVert  {u}\rVert^2_{L^{\frac{n}{2}}(\Omega)} \big) .
\end{equation}
  where $2^*=2n/(n-2)$.
\end{theorem}

\section{A minimax principle and the proof of Theorem \ref{main theorem 1}}

In this section, we first prove an adapted version of  variational principle
presented in Theorem \ref{variational} which is  applicable specifically to our
problem when \(\mu<0\), and $p>2$. Afterwards, we proceed with the proof of
Theorem \ref{main theorem 1}.

 We  consider the Banach space
\(
V= \{ u\in H_0^1(\Omega)\cap L^p(\Omega),\nabla . u=0\}
\)
equipped with the  norm
\[
\rVert u \rVert:= \rVert u \rVert_{ H_0^{1}(\Omega)}
+\rVert u \rVert_{ L^{p}(\Omega)}.
\]

Let \(\Lambda u\) be the operator \(\Lambda u := (u . \nabla)u\), that is
\[
\langle \Lambda u, v\rangle
= \int_\Omega (\Lambda u)v =\int_\Omega\sum_{j,k=1}^n
u_k\frac{\partial u_j}{\partial x_k}v_j.
\]

Let \( K\) be a convex and weakly closed subset of \(V\).
As stated in Theorem \ref{variational} we shall consider the functional
\( M : K \times K \to\mathbb{R}\) given in  \eqref{def II}.

\begin{proof}[Proof of Theorem \ref{variational}]
It follows from condition (i) in the theorem that there exists \(\bar{u} \in K\)
such that
\begin{equation}\label{ineq 1}
    \frac{1}{2}\int_{\Omega} |\nabla \bar{u}|^2\,dx
    -\frac{1}{2}\int_{\Omega} |\nabla v|^2\,dx
    \leq \int_\Omega (   f(x)+\frac{1}{p}|\bar{u}|^{p}- \Lambda \bar{u}) (\bar{u}-v)\,dx ,
    \quad \forall v\in  K.
\end{equation}
It also follows from (ii) that there exists \(\bar{v} \in K\) such that
\begin{equation}\label{ineq 2}
\int_\Omega\nabla\bar{v}.\nabla\eta\,dx=\int_\Omega (f(x)
+|\bar{u}|^{p-2}\bar{u} - \Lambda \bar{u})\eta\,dx,\quad \forall\eta\in V.
\end{equation}
Substituting \(\eta = \bar{u} - \bar{v}\) in the latter equality gives
\begin{equation}\label{ineq 3}
\int_\Omega\nabla\bar{v}.\nabla(\bar{u} - \bar{v})\,dx
=\int_\Omega (f(x)  +|\bar{u}|^{p-2}\bar{u} - \Lambda \bar{u})(\bar{u}
- \bar{v})\,dx,\quad \forall\eta\in V.
\end{equation}
Setting \(v = \bar{v}\) in \eqref{ineq 1} and taking into account the equality
\eqref{ineq 3} we obtain that
\begin{equation}\label{ineq 4}
 \int_\Omega\nabla\bar{v}.\nabla(\bar{u} - \bar{v})\,dx
 \geq\frac{1}{2}\int_{\Omega} |\nabla \bar{u}|^2\,dx
 -\frac{1}{2}\int_{\Omega} |\nabla \bar{v}|^2\,dx.
\end{equation}
On the other hand, it follows from the convexity of \( g(t)=\frac{1}{2}t^2\) that
\begin{equation}\label{ineq 5}
    \int_\Omega\nabla\bar{v}.\nabla(\bar{u} - \bar{v})\,dx \leq\frac{1}{2}\int_{\Omega} |\nabla \bar{u}|^2\,dx -\frac{1}{2}\int_{\Omega} |\nabla \bar{v}|^2\,dx.
\end{equation}
Inequalities \eqref{ineq 4} and \eqref{ineq 5} together imply that
\[
\int_\Omega\nabla\bar{v}.\nabla(\bar{u} - \bar{v})\,dx
=\frac{1}{2}\int_{\Omega} |\nabla \bar{u}|^2\,dx-\frac{1}{2}\int_{\Omega}
|\nabla \bar{v}|^2\,dx.
\]
Therefore,
\[
\int_{\Omega} |\nabla \bar{u}-\nabla \bar{v}|^2=0,
\]
from which it follows that \(\bar{v}=\bar{u}\) for a.e.\
 \(x \in\Omega\). Hence, the equality \eqref{ineq 2} proves the desired  result.
\end{proof}

We shall apply Theorem \ref{variational} to prove the existence of solution  in
Theorem \ref{main theorem 1}.  The convex subset \(K\) of \(V\)
required in Theorem \ref{variational} is defined by
\begin{equation}\label{def K}
K(r)=\{u\in V: \rVert u \rVert_{ W^{2,2}(\Omega)}\leq r\},
\end{equation}
for some \(r > 0\) to be determined. To see that \(K(r)\) is weakly closed,
we present the proof of this statement in the following lemma.

\begin{lemma}\label{Kclosed}
Let \(r > 0\) be fixed. The set
\[
K(r)=\{u\in V: \rVert u \rVert_{ W^{2,2}(\Omega)}\leq r\}
\]
is weakly closed in \(V\).
\end{lemma}

\begin{proof}
Let \(\{u_m\}\) be a sequence in \(K(r)\) such that \(u_m\rightharpoonup u\) weakly
in \(V\).
Then there exists a subsequence of \(u_m\), denoted by \(u_m\) again such that
\(u_m\to u\) a.e in \(\Omega\).
On the other hand,  \(\rVert u_m \rVert_{ W^{2,2}(\Omega)}\leq r\) for all
\(m \in \mathbb{N}\) and so  \(\{u_m\}\)
is bounded in \(W^{2,2}(\Omega)\). Going if necessary to a subsequence, there exists
\(\bar{u}  \in W^{2,2}(\Omega)\)  such that \(u_m \rightharpoonup \bar{u}\) weakly in
\(W^{2,2}(\Omega)\) and \(u_m(x) \to\bar{u}(x) \) for a.e.\
 \(x \in \Omega\). It follows then \(u(x)=\bar{u}(x) \) for a.e. \(x \in \Omega\).
Thus \(u_m \rightharpoonup u\) weakly in \(W^{2,2}(\Omega)\). Now from the
weak lower semi-continuity of the norm in \(W^{2,2}(\Omega)\)  follows that
\[
\rVert u \rVert_{ W^{2,2}(\Omega)}
\leq \liminf_{m\to\infty} \rVert u_m \rVert_{ W^{2,2}(\Omega)}\leq r ,
\]
which means that \(u\in K(r)\).
\end{proof}

To apply Theorem \ref{variational}, we need to verify both conditions (i)
and (ii) in this Theorem. To verify
condition (i) we  use the following version of the well-known Ky Fan's
min-max principle \cite{HB}. We refer to \cite[Lemma 12.1]{NG} for a proof.

\begin{lemma}\label{Ky Fan 2}
Let \(E\) be a closed convex subset of a reflexive Banach space \(H\), and consider
\(M : E \times E \to\bar{\mathbb{R}}\)
to be a functional such that:
\begin{itemize}
    \item[(1)] For each \(y\in E\), the map \(x \to M(x, y)\) is weakly lower semi-
     continuous on \(E\).
    \item[(2)] For each \(x \in E\), the map \(y \to M(x, y)\) is concave on \(E\).
    \item[(3)] There exists \( \gamma\in \mathbb{R}\) such that \(M(x, x) \leq \gamma\)
    for every  \(x \in E\).
    \item[(4)] There exists a \(y_0 \in E\) such that
     \(E_0 = \{x \in E : M(x, y_0) \leq \gamma\} \) is bounded.
\end{itemize}
Then there exits \(\bar{x} \in E\) such that \(M(\bar{x}, y) \leq \gamma\) for all
\(y \in E\).
\end{lemma}

One of the requirements in Lemma \ref{Ky Fan 2} is  the lower semi-continuity of
$M(.\;,v)$ for a fixed $v$. To verify  that, we begin with the following Lemma.

\begin{lemma}\label{lem SC 1}
For each \( v\in K\), the map   \( u\to\langle \Lambda u, v\rangle  \) is weakly
continuous on \(K\) for the values of  \( n,p \) in Theorem \ref{main theorem 1}.
\end{lemma}

\begin{proof}
Fix \( v \in  K\), and let  \(u^m\rightharpoonup u\) weakly in \(K\). We have
\begin{align*}
\big| \langle \Lambda u^m, v\rangle -  \langle \Lambda u, v\rangle  \big|
&=\Big|\sum_{j,k=1}^n\int_\Omega\big(u_k^m\frac{\partial u_j^m}{\partial x_k}v_j
 -u_k\frac{\partial u_j}{\partial x_k}v_j\big) \,dx\Big| \\
&=\Big|\sum_{j,k=1}^n\int_\Omega\big( (u_k^m-u_k)\frac{\partial u_j^m}{\partial x_k}v_j
 +u_k\frac{\partial (u_j^m-u_j)}{\partial x_k}v_j\big) \,dx\Big| \\
&\leq \sum_{j,k=1}^n\int_\Omega\Big| (u_k^m-u_k)\frac{\partial u_j^m}{\partial x_k}v_j
 \Big|
+\sum_{j,k=1}^n\int_\Omega \Big| u_k\frac{\partial (u_j^m-u_j)}{\partial x_k}v_j\Big|
\,dx.
\end{align*}
On the other hand, by H\"{o}lder inequality we conclude that
\[
\sum_{j,k=1}^n\int_\Omega\Big| (u_k^m-u_k)\frac{\partial u_j^m}{\partial x_k}v_j\Big|
    \leq \rVert (u^m-u)v \rVert_{L^2}\rVert \nabla u^m \rVert_{L^2}
    \leq \rVert u^m-u \rVert_{L^4}\rVert v \rVert_{L^4} \rVert \nabla u^m \rVert_{L^2}.
\]
Therefore,
\[
\big| \langle\Lambda u^m, v\rangle -  \langle\Lambda u, v\rangle \big|
\leq \rVert u^m-u \rVert_{L^4}\rVert v \rVert_{L^4} \rVert \nabla u^m \rVert_{L^2}
 +\sum_{j,k=1}^n\int_\Omega \Big| u_k\frac{\partial (u_j^m-u_j)}{\partial x_k}v_j\Big|
 \,dx.
\]
Moreover, since the space \(W^{2,2}(\Omega)\) is compactly imbedded into
\(L^4(\Omega)\), for all \(n\leq 7\),  it follows that \(u^m\to u\) strongly
in \(L^4(\Omega)\).
Furthermore, since $u,v$ are in \(W^{2,2}(\Omega)\), we deduce
from H\"{o}lder's inequality that \(u_kv_j\in L^2(\Omega)\).
Finally,  since  \(\nabla u^m\rightharpoonup \nabla u\) weakly in \(K\),  by
definition of weak convergence  the result follows.
\end{proof}

\begin{lemma} \label{WLSC}
Fro each  \( v\in K\), the map \(u\to M(u, v) \) is weakly lower semi-continuous on
\(K\) for the values of  \( n,p \) in Theorem \ref{main theorem 1}.
\end{lemma}

\begin{proof}
 Let \( v \in  K\) be fixed and \(u_m\rightharpoonup u\) weakly in \(K\).
Since $<\Lambda u , u> = 0 $ resulting from \eqref{def II} we have
\begin{equation}\label{M(u)}
\begin{aligned}
M(u,v)
& =\frac{1}{2}\int_{\Omega} |\nabla u|^2\,dx -\langle\Lambda u , v\rangle
 -\int_{\Omega} f(x)u\,dx-\frac{1}{p}\int_{\Omega}  |u|^{p}\,dx  \\
&\quad +  \frac{1}{p}\int_{\Omega}  |u|^{p-2}uv\,dx+ \frac{1}{2}
\int_{\Omega} |\nabla v|^2\,dx+\int_{\Omega} f(x)v\,dx,
\end{aligned}
\end{equation}
Now we shall verify lower semi-continuity of every single part in \eqref{M(u)} separately.
Note that the last two terms in \eqref{M(u)} are constant  with respect to  $u$.
\begin{itemize}
 \item Since the function \(g(u) = |u|^2 \) is convex, it can easily be shown that
\[
 \int_{\Omega}|\nabla u|^2\,dx\leq \liminf_{m \to \infty}\int_{\Omega}|\nabla u^m|^2\,dx .
\]
 that implies the map \( u\to\int_{\Omega}|\nabla u|^2  \,dx\) is weakly lower
 semi-continuous.

 \item The map \( u\to-\int_{\Omega}\Lambda u. v \,dx\) is weakly lower semi-continuous
 by Lemma \ref{lem SC 1}.

 \item  Since \(f \in L^2(\Omega)\), applying the definition of weak convergence
 leads to
 \[
 \int_{\Omega}f(x)\  u  \,dx=\liminf_{n\to \infty}\int_{\Omega}f(x) u^m\,dx.
 \]

 \item The map \( u\to\int_{\Omega}|u|^p\,dx\) is weakly lower semi-continuous for
 \( n,p \) in Theorem \ref{main theorem 1} because
\begin{itemize}
       \item[(i)] if \(n\leq 4\), then  \(W^{2,2}(\Omega)\) is compactly imbedded
       into \( L^p(\Omega)\) for all \(  p>2\), and
       \item[(ii)] if \(5\leq n\leq 7\), then \(W^{2,2}(\Omega)\) is compactly imbedded
       into \( L^p(\Omega)\) for all \( 2<p< \frac{2n}{n-4}\).
\end{itemize}

It then follows for both of cases that
\[
\lim_{n\to \infty}\int_{\Omega}| u_m|^P=\int_{\Omega}| u|^P\,dx ,
\]

\item The map \( u\to\int_{\Omega}|u|^{p-2}uv\,dx\) is weakly lower semi-continuous for
\( n,p \) in Theorem \ref{main theorem 1} because
\begin{itemize}
   \item[(i)] if \(n\leq 4\), then  \(W^{2,2}(\Omega)\) is compactly  imbedded into
   \( L^{2(p-1)}(\Omega)\) for all \(  p>2\), and
   \item[(ii)] if \(5\leq n\leq 7\), then \(W^{2,2}(\Omega)\) is compactly imbedded
   into \( L^{2(p-1)}(\Omega)\) for all \( 2<p< \frac{2n-4}{n-4}\).
\end{itemize}
In both cases, we have  $|u|^{p-2}u\in L^2(\Omega) $ from which we deduce that the map
\( u\to\int_{\Omega}|u|^{p-2}uv\,dx\) is continuous functional and
\[
\lim_{n\to \infty}
\int_{\Omega}|u_m|^{p-2}uv\,dx= \int_{\Omega}|u|^{p-2}uv\,dx.
 \]
\end{itemize}
This completes the proof.
\end{proof}

We are now in a position to state the following result addressing condition (i)
in Theorem \ref{main theorem 1}.

\begin{lemma}\label{cond i}
Let \( K=K(r)\) be a convex and weakly closed subset of \(V\) defined in \eqref{def K}.
Let \( M : K \times K \to\mathbb{R}\) be defined as \eqref{def II} and \(n,p\) as in
Theorem \ref{main theorem 1}. Then there exists \(\bar{u} \in K\)    such that
\[M(\bar{u}, v)\leq 0\quad  \forall v\in  K.\]
\end{lemma}

\begin{proof}
 We shall show that the function \(M\) satisfies all the conditions of the Ky Fan's
Min-Max Principle presented in Lemma \ref{Ky Fan 2}.
The condition (1) is provided by Lemma \ref{WLSC}. For
each \(u\in K\), the map \(v \to M(u, v)\) is concave on \(K\) since \(M(u, v)\) is a
linear functional with respect to \(v\) except
\( -\frac{1}{2}\int_{\Omega} |\nabla v|^2\,dx\), which is in fact concave.
 Also we have \( M(u, u) = 0 = \gamma\) for every \(u \in K\). Finally, since
\(u\in K\) we have that \(\rVert u \rVert_{ W^{2,2}(\Omega)}\leq r\).
Thus, we can conclude that
\( \{u \in K : M(u, v) \leq 0\}\) is bounded. It now follows by Lemma \ref{Ky Fan 2}
that there exists \(\bar{u} \in K\) such that
\[
M(\bar{u}, v)\leq 0\quad  \forall v\in  K,
\]
as desired.
\end{proof}

Our next task consists of verifying condition (ii) in Theorem \ref{main theorem 1}.
To do this, we start with the following two lemmas, which provide us the required
estimates.
Hereafter $C$ will denote a positive constant, not necessarily the same one.

\begin{lemma}\label{C1C2}
 Let \(\Omega \subset \mathbb{R}^n \) be a bounded domain and \( 1<p\).
Then for any $\ u\in K(r) $ we have
\[
\big\rVert f+|u|^{p-1}u-\Lambda u\big\rVert_{L^2(\Omega)}
\leq \big( \rVert f\rVert_{L^2(\Omega)}
+\rVert u\rVert^{p-1}_{L^{2(p-1)}}+\rVert {u}\rVert^{2}_{W^{1,2^*}(\Omega)}
\rVert  {u}\rVert^2_{L^{\frac{n}{2}}(\Omega)} \big).
\]
\end{lemma}

\begin{proof}
Let \(u \in K(r)\). By H\"older's inequality we have
\begin{align*}
     \big\rVert f+|u|^{p-2}u-\Lambda u\big\rVert_{L^2(\Omega)}
     &\leq \rVert f\rVert_{L^2(\Omega)}+\rVert u^{p-1}\rVert_{L^2(\Omega)}
     +\rVert \Lambda u\rVert_{L^2(\Omega)}\\
     &\leq \rVert f\rVert_{L^2(\Omega)}+\rVert u\rVert^{p-1}_{L^{2(p-1)}(\Omega)}
     +\rVert \nabla u\rVert^2_{L^{2^*}(\Omega)}
     \rVert  u\rVert^2_{L^{\frac{n}{2}}(\Omega)}, \quad (\text{where }2^*=\frac{2n}{n-2})\\
     &\leq \rVert f\rVert_{L^2(\Omega)}+\rVert u\rVert^{p-1}_{L^{2(p-1)}(\Omega)}
     +\rVert  u\rVert^2_{W^{1,2^*}(\Omega)}\rVert  u\rVert^2_{L^{\frac{n}{2}}(\Omega)}.
\end{align*}
as desired.
\end{proof}

\begin{lemma}\label{rbar}
Let \( p>2 \) and \(C>0\) be given. Then there exists \(0<{r}\in \mathbb{R}\) which satisfies
 \[ C (\rVert f\rVert_{L^2(\Omega)} +{r}^{p-1}+{r}^4)\leq {r},\]
 where \(\rVert f\rVert_{L^2(\Omega)}\) be small enough.
\end{lemma}

\begin{proof} Since \(p >2\), we can choose \({r}\) such that
\[
C( {r}^{p-1}+ {r}^4)\leq \frac{{r}}{2}.
\]
Now if \(C\rVert f\rVert_{L^2(\Omega)}\leq \frac{{r}}{2}\) then we have
\[
C (\rVert f\rVert_{L^2(\Omega)} +{r}^{p-1}+{r}^4)\leq {r},
\]
as desired.
\end{proof}

Here is another useful results that we shall use in the sequel.
See \cite[Theorem 1.2]{RH} for a more general version of the following result.

\begin{lemma}\label{nabla lem}
If \(g\in L^2(\Omega)\), then there exists  \(u \in W^{2,2}(\Omega)\cap H_0^1(\Omega)\),
a scalar function \(P : \Omega \to\mathbb{R}\)  and a constant \( C\) such that
 \[
\Delta u + \nabla P = g, \quad  \nabla\cdot u=0,\quad  u|_{\partial \Omega}=0,
\]
and
\[\
rVert \Delta u \rVert_{L^2(\Omega)}+\rVert \nabla P \rVert_{L^2(\Omega)}
\leq C\rVert g \rVert_{L^2(\Omega)}.
\]
\end{lemma}

The following inequality is proved in \cite[Lemma 9.17]{DN}.

\begin{lemma} \label{Lu}
Let \(\Omega \) be a bounded \( C^{1,1 }\) domain in \(\mathbb{R}^n\) and let the
operator \(Lu = a^{ij}(x)D_{ij}u+b^i(x)D_iu+c(x)u\)
 be strictly Elliptic in \(\Omega \) with coefficients
 \( a_{ij} \in C(\Omega )\), \(b_i , c \in L^\infty(\Omega) \), with
 \(i, j = 1, \dots, n\) and \(c \leq 0\). Then there
 exists a positive constant \(C\) (independent of \(u\)) such that
\[
 rVert u \rVert_{ W^{2,p}(\Omega)}\leq C \rVert Lu \rVert_{ L^{p}(\Omega)},
\]
for all \(u \in  W^{2,p}(\Omega) \cap W_{0}^{1,p}(\Omega),\  1 < p < \infty\).
\end{lemma}

Here comes a direct consequence of Lemma \ref{Lu}.

\begin{corollary} \label{cor lap}
Let \(\Omega \) be a bounded \( C^{1,1 }\) domain in \(\mathbb{R}^n\).
Then there exists  a constant \(C\) such that
\[
\rVert u \rVert_{ W^{2,2}(\Omega)}\leq C \rVert \Delta u \rVert_{ L^2(\Omega)},
\]
for all \(u \in  W^{2,2}(\Omega) \cap H_0^1(\Omega)\).
\end{corollary}

We are now ready to prove Theorem \ref{main theorem 1}.


\begin{proof}[Proof of Theorem \ref{main theorem 1}]
 Without loss of generality we may suppose that \(\mu=-1\).
We define \(K = K({r})\) for $r>0$ to be determined presently.
By Lemma \ref{cond i} we have the existence of a non-trivial \(\bar{u} \in K\)
such that
\[
M(\bar{u}, v)\leq 0\quad  \forall v\in  K.
\]
Now we shall show the existence of \(\bar{v}\) that satisfy condition (ii) in the
theorem. Consider
\[
g(x)=f+|\bar{u}|^{p-1}\bar{u}-\Lambda \bar{u}.
\]

Thus we have to show there exists \(\bar{v}\in K\)  that the following equation holds
in the weak sense,
\begin{equation}\label{eq. lap}
   -\Delta v + \nabla P=g(x).
\end{equation}
By Lemma \ref{nabla lem} there exists \(\bar{v}\in V \)
which satisfies \eqref{eq. lap} and
\begin{equation}\label{nabla}
\rVert \Delta \bar{v} \rVert_{L^2(\Omega)}
+\rVert \nabla P \rVert_{L^2(\Omega)}\leq C\rVert g \rVert_{L^2(\Omega).}
\end{equation}
It is sufficient to show that \( \bar{v}\in K\). The  estimate \eqref{nabla}
together with  Lemma \ref{C1C2} imply that
\begin{equation} \label{nabla111}
\begin{aligned}
   \rVert \Delta {\bar{v}} \rVert_{L^2(\Omega)}
   +\rVert \nabla P \rVert_{L^2(\Omega)}
   &\leq C \rVert f+|\bar{u}|^{p-1}\bar{u}-\Lambda \bar{u}\rVert \\
   &\leq C\big( \rVert f\rVert_{L^2(\Omega)}+\rVert \bar{u}\rVert^{p-1}_{L^{2(p-1)}}
   +\rVert \bar{u}\rVert^{2}_{W^{1,2^*}(\Omega)} \rVert
    \bar{u}\rVert^2_{L^{\frac{n}{2}}(\Omega)} \big) .
\end{aligned}
\end{equation}

On the other hand, Corollary \ref{cor lap} together with \eqref{nabla111} yield that
\begin{equation}\label{CC}
\begin{aligned}
    \rVert \bar{v} \rVert_{ W^{2,2}(\Omega)}
     &\leq C \rVert \Delta \bar{v} \rVert_{ L^2(\Omega)}
     \leq {C}\big(\rVert \Delta \bar{v} \rVert_{L^2(\Omega)}
    +\rVert \nabla P \rVert_{L^2(\Omega)} \\
     &\leq C\big( \rVert f\rVert_{L^2(\Omega)}
    +\rVert \bar{u}\rVert^{p-1}_{L^{2(p-1)}}
    +\rVert \bar{u}\rVert^{2}_{W^{1,2^*}(\Omega)} \rVert
    \bar{u}\rVert^2_{L^{\frac{n}{2}}(\Omega)} \big).
\end{aligned}
\end{equation}


From the  imbeddings of  \(W^{2,2}(\Omega)\hookrightarrow L^{2(p-1)}\) and
\(W^{2,2}(\Omega)\hookrightarrow W^{1,2^*}(\Omega)\)  we obtain from \eqref{CC}
that
\begin{equation}\label{CC2}
  \rVert \bar{v} \rVert_{ W^{2,2}(\Omega)}
  \leq C\big( \rVert f\rVert_{L^2(\Omega)}+\rVert \bar{u}\rVert^{p-1}_{w^{2,2}(\Omega)}
  +\rVert \bar{u}\rVert^{2}_{W^{2,2}(\Omega)} \rVert  \bar{u}\rVert^2_{W^{2,2}(\Omega)}
  \big).
\end{equation}
Let $r$ be as in Lemma \ref{rbar} for $C$ given in the last inequality above.
The inequality \eqref{CC2} and Lemmas \ref{rbar} yield that
\[
  \rVert \bar{v} \rVert_{ W^{2,2}(\Omega)}
\leq   C(\rVert f\rVert_{L^2(\Omega)} +{r}^{p-1}+{r}^4)\leq {r},
\]
where \(\rVert f\rVert_{L^2(\Omega)}\) is small enough.
That means \( \bar{v}\in K\) and so $\bar{v}=\bar{u}$.
This completes the proof of (i) and (ii).
Now  the inequality \eqref{nabla111}  concludes that
  \begin{equation} 
  \rVert \Delta {u} \rVert_{L^2(\Omega)}+\rVert \nabla P \rVert_{L^2(\Omega)}
  \leq C\big( \rVert f\rVert_{L^2(\Omega)}+\rVert u\rVert^{p-1}_{L^{2(p-1)}}
  +\rVert {u}\rVert^{2}_{W^{1,2^*}(\Omega)}
  \rVert  {u}\rVert^2_{L^{\frac{n}{2}}(\Omega)} \big) .
\end{equation}
\end{proof}

\begin{thebibliography}{99}
\bibitem[1]{HB} H. Brezis, L. Nirenberg, G. Stampacchia; \emph{A remark on Ky Fan's Minimax Principle}, Bollettino U. M. I (1972), 293-300
\bibitem[2]{XC} X. Cai and Q. Jiu; \emph{Weak and strong solutions for the incompressible Navier-Stokes equations with damping}, J. Math. Anal. Appl., 343 (2008), 799-809.
\bibitem[3]{RH} R. Farwig and H. Sohr; \emph{Generalized resolvent estimates for the Stokes system in bounded and unbounded domains}, J. Math. Soc. Japan, vol 46 no 4 (1994), 607-643.
\bibitem[4]{DN} D. Gilbarg, N. S. Trudinger; \emph{Elliptic partial differential equations of second order}, Reprint of the (1998) edition. Classics in Mathematics. Springer-Verlag, Berlin, (2001).
\bibitem[5]{NG} N. Ghoussoub; \emph{Self-dual Partial Differential Systems and Their Variational Principles}, Springer Monographs in Mathematics, Springer, New York (2008).
\bibitem[6]{EH} E. Hopf; \emph{Über die Anfangswertaufgabe für die hydrodynamischen Grundgleichungen}, Math. Nachr. 4 (1951), 213-231.
\bibitem[7]{LHS} L. Hsiao; emph{Quasilinear Hyperbolic Systems and Dissipative Mechanisms}, World Scientific, 1997.
\bibitem[8]{YJX} Y. Jia, X. W. Zhang, B. Q. Dong; \emph{The asymptotic behavior of solutions to three dimensional Navier-Stokes equations with nonlinear damping}, Nonlinear Anal. Real World Appl., 12 (2011), 1736-1747.
\bibitem[9]{ZHJ} Z. H. Jiang; \emph{Asymptotic behavior of strong solutions to the 3D Navier-Stokes equations with a nonlinear damping term}, Nonlinear Anal., 75 (2012), 5002-5009.
\bibitem[10]{ZHM} Z. H. Jiang, M. X. Zhu; \emph{The large time behavior of solutions to 3D Navier-Stokes equations with nonlinear damping}, Math. Methods Appl. Sci., 35 (2012), 97-102.
\bibitem[11]{JL} J. Leray; \emph{Sur le mouvement d'un liquide visqueux emplissant l'espace}, Acta Math. 63 (1934), 193-248.
\bibitem[12]{ZLD} Z. Li, D. Shi, M. Li; \emph{Stabilized mixed finite element methods for the Navier-Stokes equations with damping}, Math. Methods Appl. Sci. 42 (2019), 605-619.
\bibitem[13]{WLW} W. Li, W. Wang, Q. Jiu; \emph{Existence and uniqueness of the weak solutions for the steady incompressible Navier Stokes equations with damping}, African diaspora journal of mathematics, volume 12, number 2 (2011), 57-72.
\bibitem[14]{LML} M. Li, Z. Li, D. SHI; \emph{Unconditional Optimal Error Estimates for the Transient Navier-Stokes Equations with Damping}. Adv. Appl. Math. Mech., v. 14, n. 1 (2022), 248-274.
\bibitem[18]{HQL} H. Qiu, L. Mei; \emph{Multi-level stabilized algorithms for the stationary incompressible Navier-Stokes equations with damping}, Appl. Numer. Math., 143 (2019), pp.188-202.
\bibitem[19]{HQY} H. Qiu, Y. Zhang, L. Mei; \emph{A mixed-FEM for Navier-Stokes type variational inequality with nonlinear damping term}, Comput. Math. Appl., 73(10) (2017), 2191-2207
\bibitem[20]{XLS} X. L. Song, F. Liang, J. H. Wu; \emph{Pullback D-attractors for three-dimensional Navier-Stokes equations with nonlinear damping}, Bound. Value Probl., (2016), Paper No. 145, 15 pp.
\bibitem[23]{RTE} R. Temam; \emph{Navier-Stokes Equations, Theory and Numerical Analysis. North-Holland Publishing Company}, Amsterdam, New York, Oxford, 1979.
\bibitem[24]{YZ} Y. Zhou; \emph{Regularity and uniqueness for the 3D incompressible Navier-Stokes equations with damping}, Appl. Math. Lett., 25 (2012), 1822-1825.
\bibitem[25]{ZJZ} Z. J. Zhang, X. L. Wu, M. Lu; \emph{On the uniqueness of strong solution to the incompressible Navier-Stokes equations with damping}, J. Math. Anal. Appl., 377 (2011), 414-419.
\bibitem[26]{ZZC} Z. Zhang, C. P. Wu, Z. Yao; \emph{Remarks on global regularity for the 3D MHD system with damping}, Appl. Math and Com., 333: 1-7 (2018)
\bibitem[27]{ZBS} B. Zheng, Y. Shang; \emph{Two-level defect-correction stabilized algorithms for the simulation of 2D/3D steady Navier-Stokes equations with damping}. Applied Numerical Mathematics, v. 163 (2021) 182-203.
\end{thebibliography}
\end{document}
